\chapter{2001年上海卷（秋理）}
\section{填空题}
\begin{problem}
设函数 \( f(x) = \begin{cases} 2^{-x}, & x \in (-\infty, 1] \\ \log_{81} x, & x \in (1, +\infty) \end{cases} \)，则满足 \( f(x) = \frac{1}{4} \) 的 \( x \) 值为\fillinblank{}.
\end{problem}

\begin{problem}
设数列 \(\{a_{n}\}\) 的通项为 \(a_{n} = 2n - 7 (n \in \N)\)，则 \(|a_{1}| + |a_{2}| + \dots + |a_{10}| = \)\fillinblank{}.
\end{problem}

\begin{problem}
设 \(P\) 为双曲线 \(\frac{x^2}{4} - y^2 = 1\) 上一动点，\(O\) 为坐标原点，\(M\) 为线段 \(OP\) 的中点，则点 \(M\) 的轨迹方程为\fillinblank{}.
\end{problem}

\begin{problem}
设集合 \(A = \{x \mid 2 \lg x = \lg (8x - 15), x \in \R\}\)，\(B = \{x \mid \cos \frac{x}{2} > 0, x \in \R\}\)，则 \(A \cap B\) 的元素个数为\fillinblank{} 个.
\end{problem}

\begin{problem}
抛物线 \(x^{2} - 4y - 3 = 0\) 的焦点坐标为\fillinblank{}.
\end{problem}

\begin{problem}
设数列 \(\{a_{n}\}\) 是公比 \(q > 0\) 的等比数列，\(S_{n}\) 是它的前 \(n\) 项和. \(\lim\limits_{n \to \infty} S_{n} = 7\)，则此数列的首项 \(a_{1}\) 的取值范围是\fillinblank{}.
\end{problem}

\begin{problem}
某餐厅供应客饭，每位顾客可以在餐厅提供的菜肴中任选 \(2\) 荤 \(2\) 素共 \(4\) 种不同的品种. 现在餐厅准备了 \(5\) 种不同的荤菜，若要保证每位顾客有 \(200\) 种以上不同的选择，则餐厅至少还需要准备不同的素菜品种\fillinblank{} 种. （结果用数值表示）
\end{problem}

\begin{problem}
在代数式 \((4x^{2} - 2x - 5)\left(1 + \frac{1}{x^{2}}\right)^{5}\) 的展开式中，常数项为\fillinblank{}.
\end{problem}

\begin{problem}
设 \(x = \sin \alpha\)，\(\alpha \in \left[-\frac{\pi}{6}, \frac{5\pi}{6}\right]\)，则 \(\arccos x\) 的取值范围为\fillinblank{}.
\end{problem}

\begin{problem}
直线 \( y = 2x - \frac{1}{2} \) 与曲线 \( \begin{cases} x = \sin \varphi, \\ y = \cos 2\varphi. \end{cases} \)（\(\varphi\) 为参数）的交点坐标为\fillinblank{}.
\end{problem}

\begin{problem}
已知两个圆：\(x^{2} + y^{2} = 1\) \circled{1}与 \(x^{2} + (y - 3)^{2} = 1\) \circled{2}，则由\circled{1}式减去\circled{2}式可得上述两圆的对称轴方程. 将上述命题在曲线的情况下加以推广，即要求得到一个更一般的命题，而已知命题应成为所推广命题的一个特例，推广的命题为\fillinblank{}.
\end{problem}

\begin{problem}
据报道，我国目前已成为世界上受荒漠化危害最严重的国家之一. 下左图表示我国土地沙化总面积在上个世纪五六十年代，七八十年代，九十年代的变化情况. 由图中的相关信息，可将上述有关年代中，我国年平均土地沙化面积在右下图中图示为：

\bitmapfigure[max width=0.64\linewidth]{img/2001/shanghai_science/q12_fig1.png}
\end{problem}

\section{单选题}
\begin{problem}
\(a = 3\) 是直线 \(ax + 2y + 3a = 0\) 和直线 \(3x + (a - 1)y = a - 7\) 平行且不重合的（\quad）

\choices
    {充分非必要条件}
    {必要非充分条件}
    {充要条件}
    {既非充分也非必要条件}
\end{problem}

\begin{problem}
在平行六面体 \(ABCD - A_{1}B_{1}C_{1}D_{1}\) 中，\(M\) 为 \(AC\) 与 \(BD\) 的交点，若 \(\overrightarrow{A_1B_1} = \bs{a}\)，\(\overrightarrow{A_1D_1} = \bs{b}\)，\(\overrightarrow{A_1A} = \bs{c}\)，则下列向量中与 \(\overrightarrow{B_1M}\) 相等的向量是（\quad）

\choices
    {\(-\frac{1}{2}\bs{a} + \frac{1}{2}\bs{b} + \bs{c}\)}
    {\(\frac{1}{2}\bs{a} + \frac{1}{2}\bs{b} + \bs{c}\)}
    {\(\frac{1}{2}\bs{a} - \frac{1}{2}\bs{b} + \bs{c}\)}
    {\(-\frac{1}{2}\bs{a} - \frac{1}{2}\bs{b} + \bs{c}\)}
\end{problem}

\begin{problem}
已知 \(a\)，\(b\) 为两条不同的直线，\(\alpha\)，\(\beta\) 为两个不同的平面，且 \(a \perp \alpha\)，\(b \perp \beta\)，则下列命题中的假命题是（\quad）

\choices
    {若 \(a \parallel b\)，则 \(\alpha \parallel \beta\)}
    {若 \(\alpha \perp \beta\)，则 \(a \perp b\)}
    {若 \(a\)，\(b\) 相交，则 \(\alpha\)，\(\beta\) 相交}
    {若 \(\alpha\)，\(\beta\) 相交，则 \(a\)，\(b\) 相交}
\end{problem}

\begin{problem}
用计算器验算函数 \( y = \frac{\lg x}{x} \)（\(x > 1\)）的若干个值，可以猜想下列命题中的真命题只能是（\quad）

\choices
    {\(y = \frac{\lg x}{x}\) 在 \((1, +\infty)\) 上是单调减函数}
    {\(y = \frac{\lg x}{x}\)，\(x \in (1, +\infty)\) 的值域为 \( \left(0, \frac{\lg 3}{3}\right] \)}
    {\(y = \frac{\lg x}{x}\)，\(x \in (1, +\infty)\) 有最小值}
    {\(\lim\limits_{n\to +\infty}\frac{\lg n}{n} = 0\)，\(n \in \N\)}
\end{problem}

\section{解答题}
\begin{problem}
已知 \(a\)，\(b\)，\(c\) 是 \(\triangle ABC\) 中 \(\angle A\)，\(\angle B\)，\(\angle C\) 的对边，\(S\) 是 \(\triangle ABC\) 的面积. 若 \(a = 4\)，\(b = 5\)，\(S = 5\sqrt{3}\)，求 \(c\) 的长度.
\end{problem}

\begin{problem}
设 \(F_{1}\)，\(F_{2}\) 为椭圆 \(\frac{x^{2}}{9} + \frac{y^{2}}{4} = 1\) 的两个焦点，\(P\) 为椭圆上的一点. 已知 \(P\)，\(F_{1}\)，\(F_{2}\) 是一个直角三角形的三个顶点，且 \(|PF_{1}| > |PF_{2}|\). 求 \(\frac{|PF_{1}|}{|PF_{2}|}\) 的值.
\end{problem}

\begin{problem}
在棱长为 \(a\) 的正方体 \(OABC - O'A'B'C'\) 中，\(E\)，\(F\) 分别是棱 \(AB\)，\(BC\) 上的动点，且 \(AE = BF\).

\begin{enumerate}
\item 求证：\(A'F \perp C'E\)；
\item 当三棱锥 \(B' - BEF\) 的体积取得最大值时，求二面角 \(B' - EF - B\) 的大小. （结果用反三角函数表示）
\end{enumerate}
\end{problem}

\begin{problem}
对任意一个非零复数 \(z\) 定义集合 \(M_z = \{\omega \mid \omega = z^{2n - 1}, n \in \N\}\).

\begin{enumerate}
\item 设 \(a\) 是方程 \(x + \frac{1}{x} = \sqrt{2}\) 的一个根，试用列举法表示集合 \(M_a\). 若在 \(M_a\) 中任取两个数，求其和为零的概率 \(P\)；
\item 设集合 \(M_z\) 中只有 \(3\) 个元素，试写出满足条件的一个 \(z\) 的值，并说明理由.
\end{enumerate}
\end{problem}

\begin{problem}
用水清洗一堆蔬菜上残留的农药，对用一定量的水清洗一次的效果作如下假定：用 \(1\) 个单位量的水可洗掉蔬菜上残留农药用量的 \(\frac{1}{2}\)，用水越多洗掉的农药量也越多，但总还有农药残留在蔬菜上. 设用 \(x\) 单位量的水清洗一次以后，蔬菜上残留的农药与本次清洗前残留有农药量之比为函数 \(f(x)\).

\begin{enumerate}
\item 试规定 \(f(0)\) 的值，并解释其实际意义；
\item 试根据假定写出函数 \(f(x)\) 应该满足的条件和具有的性质；
\item 设 \(f(x) = \frac{1}{1 + x^2}\)，现有 \(a\)（\(a > 0\)）单位量的水，可以清洗一次，也可以把水平均分成 \(2\) 份后清洗两次. 试问用哪种方案清洗后蔬菜上的农药量比较少? 说明理由.
\end{enumerate}
\end{problem}

\begin{problem}
对任意函数 \(f(x)\)，\(x \in D\)，可按图示构造一个数列发生器，其工作原理如下：

\begin{circlelist}
\item 输入数据 \(x_0 \in D\)，经数列发生器输出 \(x_1 = f(x_0)\)；
\item 若 \(x_1 \notin D\)，则数列发生器结束工作；若 \(x_1 \in D\)，则将 \(x_1\) 反馈回输入端，再输出 \(x_2 = f(x_1)\)，并依此规律继续下去. 现定义 \(f(x)=\frac{4x-2}{x+1}\).
\end{circlelist}

\begin{flushright}
\end{flushright}

\begin{enumerate}
\item 若输出 \(x_0=\frac{49}{65}\)，则由数列发生器产生数列 \(\{x_n\}\). 请写出数列 \(\{x_n\}\) 的所有项；
\item 若要数列发生器产生一个无穷的常数数列，试求输出的初始数据 \(x_0\) 的值；
\item 是否存在 \(x_0\)，在输入数据 \(x_0\) 时，该数列发生器产生一个各项均为负数的无穷数列？若存在，求出 \(x_0\) 的值；若不存在，请说明理由.
\end{enumerate}

\bitmapinclude{img/2001/shanghai_science/q23_fig1.png}
\end{problem}
